%%% FIELD observational-rate-bracket-proof
Fix \(n\) at least as large as the maximum of the thresholds in \(\text{thm:known-propensity-benchmark}\) and \(\text{thm:computable-upper-bound}\).  First compare the two experiments.  If \(\phi_n:O^n\to[0,1]\) is any measurable observed-data test, define
\[
 \psi_n:O^n\times\mathcal E_{\alpha,L,\epsilon}\longrightarrow[0,1],
 \qquad
 \psi_n(o^n,e_0)=\phi_n(o^n).
\]
This map is measurable because it is the composition of \(\phi_n\) with the measurable coordinate projection.  For every \(P\in\mathcal P^{\mathrm{obs}}\), including every law in either the null class \(V(P)=0\) or an alternative class \(V(P)\ge\rho^2\), it satisfies
\[
 \mathbb E_{P^{\otimes n}}\!\left[\psi_n(O^n,e_P)\right]
 =\mathbb E_{P^{\otimes n}}\!\left[\phi_n(O^n)\right]
\]
and
\[
 \mathbb E_{P^{\otimes n}}\!\left[1-\psi_n(O^n,e_P)\right]
 =\mathbb E_{P^{\otimes n}}\!\left[1-\phi_n(O^n)\right].
\]
Here \(e_P\in\mathcal E_{\alpha,L,\epsilon}\) by the observed-image characterization and the propensity smoothness and overlap conditions built into \(\mathcal P^{\mathrm{obs}}\).  Thus every separation \(\rho\) feasible in \(\text{def:separation-radius}\) is feasible in \(\text{def:known-propensity-separation-radius}\).  Taking infima of the two feasible sets gives
\[
 \rho_\star^{\mathrm{known}\ e}(n)\le \rho_\star(n).
\]
The lower half of \(\text{thm:known-propensity-benchmark}\) therefore yields
\[
 c_{\mathrm{rate}}r_{\mathrm{or}}(n)
 \le \rho_\star^{\mathrm{known}\ e}(n)
 \le \rho_\star(n).
\]

For the other direction, \(\text{thm:computable-upper-bound}\) supplies the single data-driven test \(\phi_n^{\mathrm{up}}\) defined by \(\text{def:corrected-energy-test}\).  Its uniform type-I error is at most \(0.05=a_0\), and its uniform type-II error over \(V(P)\ge C_{\mathrm{rate}}^2R_{\mathrm{up}}(n)^2\) is at most \(0.20=b_0\).  Hence \(C_{\mathrm{rate}}R_{\mathrm{up}}(n)\) belongs to the feasible set in \(\text{def:separation-radius}\), and
\[
 \rho_\star(n)\le C_{\mathrm{rate}}R_{\mathrm{up}}(n).
\]
Combining the last two displays proves the asserted full-class bracket.  Notice that both component theorems range over the same observed image \(\mathcal P^{\mathrm{obs}}\); no interior arm-mean condition, conditional-variance lower bound, or fitted-nuisance rate is introduced by this comparison.

It remains to determine which exact power appears in the upper endpoint.  Put
\[
 a_{\mathrm{or}}=\frac{2\gamma}{4\gamma+d},
 \qquad
 a_{\mathrm{nu}}=\frac{4\gamma s}{4\gamma s+ds+2d\gamma}.
\]
The declared conditions \(d\ge1\), \(\gamma>0\), and \(s=\alpha+\beta>0\) make both denominators strictly positive.  Direct subtraction gives
\[
 a_{\mathrm{nu}}-a_{\mathrm{or}}
 =\frac{2\gamma\{s(4\gamma+d)-2d\gamma\}}
 {(4\gamma s+ds+2d\gamma)(4\gamma+d)}.
\]
Consequently,
\[
 a_{\mathrm{nu}}\ge a_{\mathrm{or}}
 \quad\Longleftrightarrow\quad
 s\ge\frac{2d\gamma}{4\gamma+d}.
\]
Since \(n^{-a}\) is nonincreasing in \(a\) for \(n\ge1\), the definition
\(R_{\mathrm{up}}=r_{\mathrm{or}}\vee r_{\mathrm{nu}}\) implies
\[
 R_{\mathrm{up}}(n)=
 \begin{cases}
 r_{\mathrm{or}}(n),&s\ge 2d\gamma/(4\gamma+d),\\
 r_{\mathrm{nu}}(n),&s<2d\gamma/(4\gamma+d).
 \end{cases}
\]
In the first region the bracket becomes
\[
 c_{\mathrm{rate}}r_{\mathrm{or}}(n)
 \le \rho_\star(n)
 \le C_{\mathrm{rate}}r_{\mathrm{or}}(n),
\]
where the positive finite constants do not depend on \(n\); hence \(\rho_\star(n)\asymp r_{\mathrm{or}}(n)\).  At the boundary \(s=2d\gamma/(4\gamma+d)\), the subtraction formula gives \(a_{\mathrm{nu}}=a_{\mathrm{or}}\), so the definitions give the exact identity \(r_{\mathrm{nu}}(n)=r_{\mathrm{or}}(n)\) for every \(n\).  The attaining upper bound therefore contains no logarithmic multiplier at equality.  When \(s<2d\gamma/(4\gamma+d)\), the proof gives only the displayed bracket with upper endpoint \(r_{\mathrm{nu}}\), exactly as claimed.
